% English draft based on appendices/appendix_c_part1_mapping_ja.tex
\section{Mapping to Intelligence Part I}

\noindent\textbf{Scope of the account.} This appendix records editorial proposals for Part I, not their adoption. The proposed additions must retain the distinction between conditional formal statements, cognitive interpretations, and open mappings. References here to proof, identity, or refinement do not certify those stronger readings or resolve the four-layer inheritance question.

The three-layer profiles and generation or necessity wording in this appendix remain preserved, unimplemented Part I proposals. They are not proof results or revision instructions adopted by the active support-role account in Section 5. Its scoped roles and analytical boundaries do not settle these historical profile assignments or authorize changing Part I.


Part II proposes refinements to several claims from Intelligence Part I; these editorial proposals have not been adopted here.

\subsection{Individual Non-Closure}

Part I stated that individual intelligence does not close. Part II preserves a historical local-state model's conditional formal statement: the inference layer $L_F$ has no internal fixed point under fixed $\theta$. This refines non-closure into the non-closure theorem.

This remains an unimplemented editorial proposal for Part I, not a derivation from the active record/position contract or a proof of rule-changing capacity.

\subsection{Branching Explosion}

Part I described branching explosion as the proliferation of conceptual branches. Part II proposes a connection between branching explosion and the absence of an internal stopping condition combined with generative inference.

\subsection{Physarum and Transformer}

Part I treated Physarum and Transformer as cross-domain instances of intelligence. Part II re-describes them as different three-layer profiles. Physarum externalizes logs into the environment while generating environmental stopping. Transformer internalizes frozen logs but receives stopping externally.

\subsection{Gettier}

Part I treated the Gettier problem structurally. Part II connects it to $\mu_\theta$: a system that can move thresholds can also be wrong. Error tolerance and threshold mobility are structurally coupled.

\bigskip
